% ======================================================================
% Apuntes de Lectocomprensión Jurídica en Inglés — UBA
% Documento generado a partir de apuntes de clase (formación del contrato,
% common law), método Cornell.
% ======================================================================

\documentclass[12pt,oneside]{book}

% ---------------------------- IDIOMA Y CODIFICACIÓN ----------------------------
\usepackage[utf8]{inputenc}
\usepackage[T1]{fontenc}
\usepackage[spanish,es-nodecimaldot]{babel}

% ---------------------------- GEOMETRÍA ----------------------------
\usepackage[
    top=2.5cm,
    bottom=2.5cm,
    left=3cm,
    right=2.5cm,
    headheight=15pt
]{geometry}

% ---------------------------- MATEMÁTICA ----------------------------
\usepackage{amsmath}
\usepackage{amssymb}
\usepackage{mathtools}

% ---------------------------- IMÁGENES ----------------------------
\usepackage{graphicx}
\usepackage{float}
\usepackage{caption}
\usepackage{subcaption}

% ---------------------------- TABLAS ----------------------------
\usepackage{booktabs}
\usepackage{array}
\usepackage{longtable}
\usepackage{tabularx}

% ---------------------------- LISTAS ----------------------------
\usepackage{enumitem}

% ---------------------------- COLORES Y CAJAS ----------------------------
\usepackage{xcolor}
\usepackage[most]{tcolorbox}

% ---------------------------- ENCABEZADOS ----------------------------
\usepackage{fancyhdr}

% ---------------------------- TIPOGRAFÍA ----------------------------
\usepackage{ragged2e}

% ---------------------------- RUTAS DE IMÁGENES ----------------------------
\graphicspath{
    {./img/}
    {./graficos/}
    {./figuras/}
}

% ======================================================================
% METADATOS DEL DOCUMENTO
% ======================================================================
\newcommand{\universidad}{Universidad de Buenos Aires (UBA)}
\newcommand{\facultad}{Facultad de Derecho}
\newcommand{\carrera}{Carrera de Abogacía}
\newcommand{\materia}{Lectocomprensión Jurídica en Inglés}
\newcommand{\titulo}{Derecho Contractual Comparado: la Formación del Contrato}
\newcommand{\autor}{Isaac Edgar Camacho Ocampo}
\newcommand{\anio}{2026}

% ---------------------------- HIPERVÍNCULOS Y METADATOS PDF ----------------------------
\usepackage[
    colorlinks=true,
    linkcolor=blue!50!black,
    citecolor=green!40!black,
    urlcolor=blue!60!black,
    pdftitle={\titulo},
    pdfauthor={\autor},
    pdfsubject={Apuntes universitarios},
    bookmarks=true
]{hyperref}

% ---------------------------- ÍNDICE ----------------------------
\setcounter{tocdepth}{2}

% ---------------------------- CONFIGURACIÓN DE PÁRRAFOS ----------------------------
\setlength{\parindent}{0pt}
\setlength{\parskip}{0.5em}

% ---------------------------- CONFIGURACIÓN DE LISTAS ----------------------------
\setlist[itemize]{
    topsep=0.3em,
    itemsep=0.2em
}
\setlist[enumerate]{
    topsep=0.3em,
    itemsep=0.2em
}

% ---------------------------- CONTROL DE VIUDAS Y HUÉRFANAS ----------------------------
\clubpenalty=10000
\widowpenalty=10000
\emergencystretch=1.5em

% ---------------------------- COMANDOS MATEMÁTICOS ----------------------------
\newcommand{\abs}[1]{\left|#1\right|}
\newcommand{\norm}[1]{\left\lVert#1\right\rVert}
\newcommand{\vect}[1]{\mathbf{#1}}
\newcommand{\deriv}[2]{\frac{d#1}{d#2}}
\newcommand{\pderiv}[2]{\frac{\partial#1}{\partial#2}}

% ======================================================================
% CAJAS ACADÉMICAS (entornos semánticos)
% ======================================================================
\newtcolorbox{definition}[1]{
    enhanced,
    breakable,
    title={Definición: #1},
    fonttitle=\bfseries,
    colback=blue!3,
    colframe=blue!50!black,
    boxrule=0.8pt,
    arc=2mm
}

\newtcolorbox{important}{
    enhanced,
    breakable,
    title={Importante},
    fonttitle=\bfseries,
    colback=yellow!8,
    colframe=orange!70!black,
    boxrule=0.8pt,
    arc=2mm
}

\newtcolorbox{example}[1]{
    enhanced,
    breakable,
    title={Ejemplo: #1},
    fonttitle=\bfseries,
    colback=green!4,
    colframe=green!40!black,
    boxrule=0.8pt,
    arc=2mm
}

\newtcolorbox{formula}[1]{
    enhanced,
    breakable,
    title={Fórmula / Regla: #1},
    fonttitle=\bfseries,
    colback=purple!4,
    colframe=purple!50!black,
    boxrule=0.8pt,
    arc=2mm
}

\newtcolorbox{exercise}[1]{
    enhanced,
    breakable,
    title={Ejercicio: #1},
    fonttitle=\bfseries,
    colback=gray!5,
    colframe=gray!60!black,
    boxrule=0.8pt,
    arc=2mm
}

\newtcolorbox{note}{
    enhanced,
    breakable,
    title={Nota},
    fonttitle=\bfseries,
    colback=white,
    colframe=gray!50,
    boxrule=0.6pt,
    arc=2mm
}

% ======================================================================
% ENCABEZADOS Y PIES DE PÁGINA
% ======================================================================
\pagestyle{fancy}
\fancyhf{}

\fancyhead[L]{\nouppercase{\leftmark}}
\fancyhead[R]{\materia}

\fancyfoot[C]{\thepage}

\renewcommand{\headrulewidth}{0.4pt}
\renewcommand{\footrulewidth}{0pt}

\fancypagestyle{plain}{
    \fancyhf{}
    \fancyfoot[C]{\thepage}
    \renewcommand{\headrulewidth}{0pt}
}

% ======================================================================
% DOCUMENTO
% ======================================================================
\begin{document}

% ----------------------------------------------------------------------
% PORTADA
% ----------------------------------------------------------------------
\begin{titlepage}

\thispagestyle{empty}

\begin{center}

\vspace*{1cm}

{\large \textsc{\universidad}}

\vspace{0.2cm}

{\large \textsc{\facultad}}

\vspace{0.2cm}

{\large \carrera}

\vspace{3cm}

{\Huge \textbf{\materia}}

\vspace{1cm}

{\LARGE \titulo}

\vspace{0.8cm}

\textit{Comprender los conceptos es el primer paso para convertir el estudio en conocimiento.}

\vspace{1.2cm}

\rule{0.70\textwidth}{0.5pt}

\vspace{0.5cm}

{\large Apuntes de estudio}

\vspace{0.5cm}

\rule{0.70\textwidth}{0.5pt}

\vfill

{\large \autor}

\vspace{0.3cm}

{\large \anio}

\end{center}

\vfill

\begin{flushleft}
\textit{Este es un modesto aporte a los estudiantes de ``Lectocomprensión Jurídica en Inglés'' no pretende sustituir el material oficial de la materia.}
\end{flushleft}

\end{titlepage}

% ----------------------------------------------------------------------
% ÍNDICE
% ----------------------------------------------------------------------
\tableofcontents

% ----------------------------------------------------------------------
% CONTENIDO
% ----------------------------------------------------------------------
\chapter[Formación del contrato en el \textit{common law}]{Derecho Contractual Comparado: la Formación del Contrato}

\section*{Temas tratados en esta clase}
\addcontentsline{toc}{section}{Temas tratados en esta clase}

\begin{itemize}
    \item La intención de contratar (\textit{intention to create legal relations}) y sus presunciones según el tipo de vínculo —familiar/afectivo versus comercial.
    \item La oferta y la aceptación (\textit{offer and acceptance}): requisitos, vigencia temporal, contraoferta y formas de comunicación.
    \item La contraprestación (\textit{consideration}): la lógica del intercambio de valor entre las partes.
    \item La certeza de los términos (\textit{certainty of terms}) y sus consecuencias sobre la validez y la exigibilidad del contrato.
    \item Comparación permanente entre estos institutos del \textit{common law} y la teoría general del contrato del derecho argentino.
\end{itemize}

\begin{important}
\textbf{¿Para qué sirve esto en mi vida profesional?} La abogacía en un mundo interconectado exige leer, entender y negociar contratos redactados en inglés: acuerdos de distribución, franquicias, servicios internacionales o jurisprudencia comparada. Comprender cómo el \textit{common law} construye la formación del contrato —intención, oferta, aceptación, contraprestación y certeza de los términos— permite identificar qué cláusula cumple cada función, anticipar cómo un tribunal extranjero interpretaría un documento y evitar falsos amigos (como \textit{consideration}, que no es ``consideración'' sino contraprestación).

\textbf{¿Es útil para mi vida? ¿Por qué?} Sí, de manera directa: gran parte de los contratos que hoy regulan la vida cotidiana —streaming, aplicaciones, comercio electrónico, alquileres temporarios, seguros de viaje— están redactados originalmente en inglés bajo estas mismas categorías. Distinguir cuándo un intercambio de mensajes constituye una oferta vinculante, cuándo una respuesta se convierte en contraoferta, o cuándo una cláusula es tan vaga que un tribunal podría anularla, es una herramienta de lectura crítica que excede lo profesional: permite firmar —o no firmar— con criterio propio cualquier acuerdo de la vida diaria.
\end{important}

\section*{Índice temático}
\addcontentsline{toc}{section}{Índice temático}

Los cuatro bloques que siguen corresponden a los cuatro elementos que, según el texto trabajado en clase, integran la formación de un contrato en el \textit{common law}. Puede pensarse la relación entre ellos con un ejemplo cotidiano: alguien le escribe a un pintor de casas. Primero debe haber intención real de contratar (no una idea al pasar); luego una oferta concreta —``te pago \$200.000 por pintar el frente en dos días''— y una aceptación en esos mismos términos; después una contraprestación —el pago del precio a cambio del trabajo—; y finalmente los términos deben ser suficientemente ciertos —qué se pinta, con qué materiales, cuándo— porque si falta precisión, el tribunal no completará lo que las partes omitieron. Los cuatro bloques son eslabones de una misma cadena: sin cualquiera de ellos, no hay contrato exigible.

\begin{table}[H]
\centering
\begin{tabularx}{\textwidth}{>{\raggedright\arraybackslash}p{0.10\textwidth} X}
\toprule
\multicolumn{2}{l}{\textbf{Unidad 1 --- La intención de contratar}} \\
\midrule
1.1 & Intención vinculante (\textit{intention}) —acuerdo preliminar, idea vaga, compromiso real. ¿Cuándo una simple idea se convierte en intención jurídicamente vinculante? \\
1.2 & Vínculos familiares y afectivos —presunción, cónyuges, amistad, refutación. ¿Por qué se presume que en la familia no hay intención de contratar? \\
1.3 & Caso \textit{Merritt v. Merritt} (1970) —separación, hipoteca, prueba escrita. ¿Qué prueba permitió refutar la presunción entre cónyuges separados? \\
1.4 & Vínculos comerciales —presunción opuesta, negociaciones previas. ¿Por qué en los negocios se presume lo contrario que en la familia? \\
\bottomrule
\end{tabularx}
\caption{Panorama de la Unidad 1}
\end{table}

\begin{table}[H]
\centering
\begin{tabularx}{\textwidth}{>{\raggedright\arraybackslash}p{0.10\textwidth} X}
\toprule
\multicolumn{2}{l}{\textbf{Unidad 2 --- Oferta y aceptación}} \\
\midrule
2.1 & \textit{Meeting of the minds} —entendimiento mutuo, cláusulas, condiciones. ¿Qué debe compartir la mente de ambas partes para que haya acuerdo? \\
2.2 & Requisitos de la oferta —claridad, propósito, vigencia temporal. ¿Qué condiciones debe reunir una propuesta para ser una oferta válida? \\
2.3 & Terminología de las partes —\textit{offeror/offeree}, sufijos \textit{-or/-ee}. ¿Cómo identificar en inglés quién origina y quién recibe un acto jurídico? \\
2.4 & Contraoferta y comunicación —\textit{counter offer}, aceptación expresa o tácita. ¿Qué ocurre cuando la aceptación modifica los términos ofrecidos? \\
\bottomrule
\end{tabularx}
\caption{Panorama de la Unidad 2}
\end{table}

\begin{table}[H]
\centering
\begin{tabularx}{\textwidth}{>{\raggedright\arraybackslash}p{0.10\textwidth} X}
\toprule
\multicolumn{2}{l}{\textbf{Unidad 3 --- La contraprestación (\textit{consideration})}} \\
\midrule
3.1 & Intercambio de promesas —\textit{promisor}, beneficio, perjuicio. ¿Qué debe recibir y qué debe ceder cada parte para que haya contraprestación? \\
3.2 & Algo de valor —apreciación pecuniaria, valor simbólico. ¿Por qué la contraprestación debe ser susceptible de valor económico? \\
3.3 & Formas de la contraprestación —dinero, bienes, servicios, abstención. ¿La contraprestación siempre implica una suma de dinero? \\
\bottomrule
\end{tabularx}
\caption{Panorama de la Unidad 3}
\end{table}

\begin{table}[H]
\centering
\begin{tabularx}{\textwidth}{>{\raggedright\arraybackslash}p{0.10\textwidth} X}
\toprule
\multicolumn{2}{l}{\textbf{Unidad 4 --- Certeza de los términos}} \\
\midrule
4.1 & Claridad y completitud —ambigüedad, vaguedad, lagunas. ¿Qué pasa cuando el contrato no especifica sus condiciones esenciales? \\
4.2 & Nulidad total o parcial —\textit{the contract will fail}, cláusula inválida. ¿El contrato entero cae o solo la cláusula imprecisa? \\
4.3 & Interpretación judicial —\textit{four corners}, usos y costumbres, tratativas. ¿Con qué elementos cuenta el tribunal para interpretar un contrato dudoso? \\
4.4 & Validez versus exigibilidad —\textit{valid}, \textit{unenforceable}, oponibilidad. ¿Puede un contrato ser válido pero no exigible ante un tribunal? \\
\bottomrule
\end{tabularx}
\caption{Panorama de la Unidad 4}
\end{table}

\section*{Transcripción temática de la clase}
\addcontentsline{toc}{section}{Transcripción temática de la clase}

La clase retoma el desarrollo de los cuatro elementos que, según el texto de lectocomprensión trabajado, integran la formación de un contrato en el \textit{common law}: \textit{intention}, \textit{offer and acceptance}, \textit{consideration} y \textit{certainty of terms}. El docente propone desmenuzar cada uno de ellos en profundidad, contrastándolos permanentemente con la teoría general del contrato del derecho argentino. Se conserva la totalidad del contenido académico desarrollado, omitiéndose únicamente los pasajes administrativos ajenos al contenido de la materia.

\section{La intención de contratar (\textit{Intention})}

\subsection{Intención vinculante (\textit{intention})}

El docente retoma el primer elemento de la formación del contrato explicado en la clase anterior: la intención. Explica que la intención de contratar es como un acuerdo preliminar: la intención de las partes de efectivamente celebrar un contrato vinculante.

\begin{note}
\textit{Cita textual del docente:} ``Una mera propuesta o una mera idea de `che, después hacemos tal cosa' no constituye un contrato; tiene que haber una intención real de obligarse legalmente.''
\end{note}

El texto trabajado utiliza la expresión \textit{only binding relationship}, y el docente se detiene en la palabra \textit{binding}: la idea de vinculante, es decir, que surte efectos jurídicos. Señala que cuando hay un conflicto y está en duda si las partes quisieron efectivamente celebrar un contrato, el tribunal debe determinar si hubo un compromiso fehaciente o si solo fue una propuesta, una idea ``en el aire''.

\begin{example}{Una idea que no es contrato}
Si a la tarde vamos al \textit{shopping}, eso no es un contrato: es una propuesta, una idea.
\end{example}

\begin{formula}{Prueba de la intención}
El texto original en inglés señala que, ante la duda, el tribunal tomará en cuenta \textit{any evidence} (cualquier prueba o material) —por ejemplo, un correo electrónico, una invitación, un compromiso— y \textit{the nature of the relationship between the parties} (la naturaleza del vínculo entre las partes).
\end{formula}

\subsection{Vínculos familiares y afectivos}

El docente explica que, como regla general (\textit{the general rule}), en las relaciones familiares o entre miembros de una familia no se presume que exista intención de celebrar un contrato vinculante. La regla alcanza a los cónyuges, a un vínculo de amistad cercana y a los miembros de una misma familia: en todos esos casos, la presunción es que no hay intención real de contratar.

Introduce luego el conector \textit{however}, que marca una idea contraria a la anterior: sin embargo, esa presunción admite ser \textbf{refutada} —no anulada ni dejada sin efecto, sino refutada— cuando una de las partes puede demostrar, con pruebas (\textit{with evidence}), que sí existió intención de celebrar un contrato pese al vínculo cercano.

\begin{definition}{Means of evidence}
En inglés, \textit{means of evidence} equivale a lo que en la teoría contractual argentina se denomina medios de prueba, y no simplemente ``evidencia'' (falso amigo de \textit{evidence}).
\end{definition}

El texto enumera distintos elementos que el tribunal puede considerar para tener por refutada la presunción: si existió un acuerdo escrito (\textit{agreement}); si intervinieron abogados en su redacción (\textit{if lawyers were involved in the agreement}); si alguna de las partes confió en la presunción de que el acuerdo era vinculante y actuó en consecuencia (por ejemplo, realizando pagos); y \textit{the actual nature of the relationship, as opposed to form or title, among others} —es decir, la naturaleza real del vínculo, más allá de la forma o el título que se le haya dado.

\subsection{Caso \textit{Merritt v. Merritt} (1970)}

El texto ilustra la refutación de la presunción familiar con un caso inglés: \textbf{\textit{Merritt v. Merritt} (1970)}. El matrimonio Merritt era propietario de una vivienda sobre la que pesaba una hipoteca (\textit{mortgage}). Al momento de separarse, el marido y la mujer firmaron un acuerdo escrito por el cual, si la mujer continuaba pagando la hipoteca, el señor Merritt le transferiría la titularidad de la propiedad una vez cancelada por completo (\textit{Mr. Merritt could transfer the property to her once the mortgage was paid in full}).

\begin{example}{Caso Merritt v. Merritt}
Cuando la señora Merritt terminó de pagar la hipoteca, el señor Merritt sostuvo que no había tenido intención de contratar y se negó a transferir la propiedad. \textit{The court ruled that there was in fact a contract}: el tribunal resolvió que efectivamente había existido un contrato y ordenó \textit{the transfer of the home to Mrs. Merritt} —la escritura translativa de dominio a favor de la esposa.

Pese al vínculo de cónyuges (\textit{despite the family relationship}), la señora Merritt pudo aportar pruebas de que las partes sí habían tenido intención de contratar (\textit{Mrs. Merritt was able to provide evidence that the parties intended to contract}), lo que bastó para refutar la presunción.
\end{example}

El docente cierra la unidad remarcando que la presunción es que en los vínculos familiares no hay intención de contratar, salvo que se puedan presentar pruebas (\textit{evidence}) en sentido contrario.

\subsection{Vínculos comerciales}

El texto pasa luego a los \textit{commercial relationships}. El docente marca el conector \textit{conversely} (por el contrario), que introduce una idea opuesta a la de los vínculos afectivos: en un entorno comercial (\textit{business setting}) la presunción es exactamente la inversa. Se presume que las partes sí tienen intención de celebrar un contrato legalmente vinculante (\textit{in legally binding contracts}).

Trabaja especialmente la palabra \textit{bargains}, que en un registro coloquial se acerca a la idea de ``regatear'' o discutir el precio, pero en sentido jurídico refiere a todo lo previo a la celebración del contrato: las negociaciones previas, las tratativas precontractuales —por ejemplo, discutir plazos de entrega, precios y descuentos por cantidad.

\begin{note}
\textit{Cita textual del docente:} ``En un ámbito comercial se presume que todas esas tratativas o negociaciones se hacen con miras a celebrar un contrato''. El docente ejemplifica con un intercambio de correos electrónicos, el envío de un presupuesto y la posterior transferencia de una seña, que en conjunto se toman como indicios de que existió intención de contratar.
\end{note}

Aclara que esta presunción también puede refutarse, pero en sentido inverso al caso anterior: aquí se trataría de demostrar que, pese al vínculo comercial, no hubo intención de celebrar un contrato —por ejemplo, cuando en el intercambio de negociaciones previas se dejó expresamente asentado que aquello no se tendría por un contrato (\textit{if there was in the exchange the provision that it will not be considered a contract}).

A partir de esta idea, se menciona el paralelismo con las cláusulas de los contratos de trabajo que aclaran expresamente que un vínculo de locación de servicios no se tendrá por relación de dependencia, cuestión cuya legitimidad —señala el docente— puede discutirse, pero que ilustra la misma lógica: lo que las partes dejan expresamente asentado por escrito.

\begin{note}
\textit{Para reflexionar:} un alumno pregunta si el texto menciona la edad o la capacidad mental como parte de este primer elemento. El docente responde que no: en la teoría contractual argentina, la capacidad —vinculada a la edad, la capacidad mental y el discernimiento— es analizada como un elemento esencial autónomo del consentimiento. En el texto trabajado, en cambio, el primer elemento se enfoca exclusivamente en el tipo de vínculo entre las partes y en cómo se presume o no la intención, no en la capacidad de quienes contratan. Es una primera diferencia relevante en la forma en que cada sistema encara la teoría contractual.
\end{note}

\section{Oferta y aceptación (\textit{Offer and Acceptance})}

\subsection{\textit{Meeting of the minds}}

El docente introduce el segundo elemento señalando que, para que un contrato sea \textit{enforceable} (exigible, ejecutable), las partes involucradas deben alcanzar un entendimiento mutuo respecto de los términos y la aceptación.

\begin{definition}{Enforceable}
Exigible o ejecutable. Si una parte no cumple, la otra puede reclamar ante un tribunal —pero solo si están dadas las condiciones para ello; de lo contrario, no hay nada que exigir.
\end{definition}

El texto habla de \textit{the terms of the agreement} o \textit{terms and conditions}, expresión habitual en los sitios de internet. El docente marca que, en la traducción literal, muchas veces se habla de ``términos y condiciones'', pero que en la tradición jurídica argentina esos \textit{terms and conditions} son, sencillamente, las cláusulas del contrato: el plazo, el precio, la forma de cumplimiento.

\begin{note}
\textit{Cita textual del docente:} ``\textit{This shared understanding, often described as a meeting of the minds}'': es una expresión muy gráfica, dice el docente —como si las mentes de ambas partes se pusieran de acuerdo. Nosotros no hablamos de \textit{meeting of the minds} sino de acuerdo de voluntades, y para que exista ese acuerdo de voluntades tienen que concurrir dos elementos distintos: \textit{offer and acceptance}.
\end{note}

\subsection{Requisitos de la oferta}

El texto define la oferta como \textit{purposeful}: con propósito, con intencionalidad —no una manifestación casual, sino un acto en el que una parte se acerca a la otra con la intención de proponer un contrato (\textit{in which one party approaches another with the intention of forming a contract}). Agrega que \textit{an offer must clearly indicate the terms of the contract}, y que \textit{the way in which an offer is communicated can take on a variety of forms}; lo determinante es cuán claramente estén expresados todos los términos de la oferta, porque los tribunales deben evaluar si esos términos fueron suficientemente precisos.

\begin{example}{Oferta demasiado vaga}
El docente ejemplifica: si propone celebrar un contrato en el que pagará un monto por la realización de una obra, pero no expresa el monto, no aclara el plazo ni especifica concretamente qué es lo que quiere que se haga, esa propuesta resulta demasiado vaga como para constituir una oferta en sentido estricto.
\end{example}

\subsection{Terminología de las partes (\textit{offeror / offeree})}

El docente se detiene en la terminología de las partes: quien realiza la oferta es el \textbf{\textit{offeror}} y quien la recibe es el \textbf{\textit{offeree}}. Propone a los alumnos identificar otras palabras del vocabulario jurídico inglés con esta misma raíz: \textit{employer/employee} (empleador/empleado), \textit{payer/payee} (quien paga/quien recibe el pago), \textit{franchisor/franchisee} (franquiciante/franquiciado), \textit{grantor/grantee} (otorgante, para el caso de poderes u otros instrumentos).

\begin{definition}{Sufijos -or / -ee}
El docente comparte su propia estrategia mnemotécnica: piensa la terminación \textit{-or} como el sujeto activo (quien origina el negocio jurídico) y la terminación \textit{-ee} como el sujeto pasivo o beneficiario (quien recibe), aunque aclara que la equivalencia no siempre es técnicamente exacta —por ejemplo, en locador/locatario, es el locatario quien paga por el uso de la cosa. Comenta además que la terminación en doble e proviene del francés, donde funcionaría originalmente como una forma participial.
\end{definition}

Finalmente, señala que en español no existe un término equivalente y natural a \textit{offeree}: la persona que recibe la oferta puede, si la acepta, convertirse en aceptante, pero no contamos con una palabra específica para el mero receptor de la oferta antes de que la acepte.

\subsection{Contraoferta y comunicación de la aceptación}

El texto retoma el ejemplo trabajado en una clase anterior sobre un collar (\textit{necklace}) para explicar la vigencia temporal de las ofertas: dice que \textit{offers can expire after a reasonable time period} —las ofertas pueden caducar tras un período de tiempo razonable, aunque no se haya fijado un plazo expreso.

\begin{example}{El collar reclamado 50 años después}
Si alguien ofrece una recompensa por encontrar un collar perdido (\textit{a reward for finding the necklace}) y una persona se presenta cincuenta años después a reclamarla (\textit{50 years after the reward}), sería ilógico considerar que la oferta seguía vigente: \textit{for the acceptance to be valid, the offer must still be open} —la oferta ya no estaba disponible para ser aceptada (\textit{off the table}).
\end{example}

El texto agrega un segundo requisito de validez para la aceptación: debe darse en los mismos términos que la oferta, porque de lo contrario constituye una \textbf{contraoferta} (\textit{counter offer}).

\begin{note}
\textit{Cita textual del docente:} el docente ejemplifica con humor: si alguien ofrece milanesas con papas fritas y la otra parte responde ``sí, pero yo no quiero milanesa ni papas fritas, quiero ensalada'', esa respuesta no está aceptando la oferta tal cual fue formulada, sino aceptando solo una parte y proponiendo nuevas condiciones: es una \textit{counter offer}.
\end{note}

Al plantearse una contraoferta, los roles se invierten: quien era \textit{offeror} pasa a ser \textit{offeree} de esa nueva propuesta, y viceversa; el proceso puede repetirse varias veces (``ir y venir'') hasta que exista una aceptación lisa y llana de una oferta concreta. Mientras haya modificaciones, en principio no hay contrato todavía.

Por último, el texto exige que la aceptación sea comunicada de manera expresa: \textit{communicated expressly, that is to say in words or in writing} (es decir, de manera oral o escrita), \textit{or a person acting on that person's behalf} —o por una persona que actúe en su representación, como un apoderado o mandatario. El docente aclara que la comunicación también puede ser implícita, a través de conductas o signos inequívocos que permitan entender que se está prestando la aceptación.

\begin{note}
\textit{Para reflexionar:} el docente compara este esquema de \textit{offer and acceptance} con la teoría del acuerdo de voluntades del derecho argentino, que exige negociaciones previas y una clara manifestación de voluntad, pero que no lo estructura formalmente como ``una parte que ofrece y otra que acepta''. Introduce también, a modo de ejemplo cotidiano, los contratos de adhesión: alquileres, servicios de internet y telefonía celular, en los que una de las partes tiene escasa o nula capacidad de negociar sus propias condiciones —a diferencia de la libertad de contratación clásica.
\end{note}

\section{La contraprestación (\textit{Consideration})}

\subsection{Intercambio de promesas}

El docente recuerda que \textit{consideration} no significa ``consideración'' —un falso amigo frecuente— sino \textbf{contraprestación}. El texto lo define como \textit{an exchange of promises}, un intercambio de promesas o compromisos, \textit{in which each party promises something to the other; in that sense, each party is a promisor}: cada parte es, a la vez, quien promete y quien recibe una promesa, de modo que los roles son intercambiables.

\begin{definition}{Consideration}
\textit{What each person promises the other is consideration}: en ese intercambio de prestaciones recíprocas, la \textit{consideration} es, precisamente, aquello que se da a cambio.
\end{definition}

\subsection{Algo de valor (\textit{something of value})}

El texto explica que la contraprestación implica un \textit{benefit} en sentido amplio para quien la recibe, \textit{or the person to whom the promise is made has to experience some detriment, that is to say, a loss or disadvantage}. El docente lo ejemplifica: si contrata a alguien para pintar su casa, su beneficio es que la casa quedará pintada; quien pinta, por su parte, sufre un \textit{detriment} en sentido amplio —no necesariamente un perjuicio, sino que debe dar algo a cambio (su tiempo, su trabajo).

El texto agrega que \textit{consideration must involve something of value}: algo de valor. El docente traduce esta expresión a la teoría argentina como \textbf{algo susceptible de apreciación pecuniaria} —es decir, que pueda expresarse en dinero, aunque también tenga un valor sentimental o simbólico que por sí solo no bastaría (por ejemplo, un recuerdo de un familiar tiene valor afectivo, pero también se le puede asignar un valor pecuniario a los fines de una eventual indemnización, del mismo modo en que se cuantifica el daño moral).

\begin{example}{Un caramelo no alcanza}
El docente aclara que no cualquier valor mínimo satisface este requisito: un caramelo, dice a modo de ejemplo, no sería representativo de un intercambio con entidad suficiente para sostener la celebración de un contrato.
\end{example}

El texto agrega que lo que \textit{induces} —motiva o lleva— a las partes a contratar es justamente el deseo de obtener aquello que la otra parte ofrece a cambio. \textit{Generally, a contract requires consideration on both sides; no transfer is required, because the concept of consideration involves negotiation}: la contraprestación no exige necesariamente una transferencia de dinero, pero sí una negociación en la que ambas partes den algo a cambio. El docente lo resume con una fórmula que utiliza también el derecho argentino: nadie celebra un contrato por nada.

\subsection{Formas de la contraprestación}

El texto aclara que la contraprestación no siempre consiste en dinero: puede tratarse de \textit{an exchange of favors} (favores o la realización de una determinada conducta), \textit{items} (bienes u objetos) o incluso \textit{to refrain from doing something} —abstenerse de realizar una determinada conducta.

\begin{example}{La parcela del cementerio familiar}
El texto ilustra la contraprestación por abstención con un caso relacionado con una parcela de un cementerio familiar: alguien se comprometió a no ocupar una determinada parcela para que la conservaran los demás miembros de la familia. El tribunal entendió que ese compromiso de no hacer —sin que mediara pago de dinero— constituía igualmente una contraprestación válida y que el acuerdo era vinculante.
\end{example}

El docente cierra la unidad preguntando cómo traducirían \textit{consideration}, y confirma junto con los alumnos: contraprestación, o prestaciones recíprocas —nunca ``consideración''.

\section{Certeza de los términos (\textit{Certainty of Terms})}

\subsection{Claridad y completitud}

El docente presenta el cuarto y último elemento: \textit{certainty of terms}. Señala que se vincula con lo ya visto sobre la oferta y la contraoferta, pero que va más allá: \textit{for a contract to be enforceable}, válido y exigible, \textit{its terms must be clear and complete} —los términos deben ser lo suficientemente claros y abarcativos como para que ambas partes comprendan exactamente qué están contratando. Si faltan condiciones, cualquiera de las partes puede luego cuestionar ante un tribunal si efectivamente eso fue lo pactado.

\begin{formula}{Ambigüedad en los términos esenciales}
\textit{Contracts may be [void/unenforceable] if key information in the terms is ambiguous} —es decir, si hay ambigüedad o vaguedad en la información esencial del contrato.
\end{formula}

\begin{example}{El objeto indeterminado del contrato}
El docente ejemplifica: vender un bote o un auto sin especificar de cuál se trata —ni color, ni modelo, ni patente— deja indeterminado el objeto del contrato. Menciona también, a modo de digresión, un programa de telerrealidad en el que se subastan contenedores de almacenamiento abandonados sin conocer su contenido: la clase coincide en que, en esos casos, el precio se paga por algo absolutamente indeterminado en cuanto a cantidad y valor.
\end{example}

El texto agrega: \textit{contracts may be incomplete if they fail to specify obligations under the contract, or if important information is missing from the terms} —las obligaciones de cada parte deben estar bien explicitadas; si falta información esencial (por ejemplo, cantidad, plazo de entrega o precio en una compraventa de mercadería), el contrato puede considerarse incompleto.

\subsection{Nulidad total o parcial}

\begin{formula}{Nulidad total o parcial}
\textit{Contracts may not be enforceable if the terms are considered too vague} —el tribunal, frente a términos vagos, tenderá a favorecer la validez del contrato y a interpretarlo de forma que sobreviva, salvo que detecte abuso de alguna de las partes.

\textit{Courts will not create entire contracts for the parties: in situations where the court cannot determine the terms of the contract, the contract will fail either in whole or in part} —si el tribunal no cuenta con información suficiente para interpretar el contrato, no lo redactará ni lo completará por las partes; en tal caso, el contrato queda sin efecto, total o parcialmente (nulidad total o parcial).
\end{formula}

\begin{note}
\textit{Cita textual del docente:} ``El tribunal no te va a resolver el contrato, te lo va a redactar —y eso no lo va a hacer. Si no está, no está.''
\end{note}

\begin{definition}{Unenforceable vs. invalid}
\textit{Unenforceable}, en este contexto, refiere a que la cláusula o estipulación imprecisa queda sin efecto (nulidad parcial), mientras que \textit{invalid} refiere a la nulidad del contrato en su totalidad.
\end{definition}

\subsection{Interpretación judicial: los usos, las tratativas y ``\textit{the four corners}''}

El texto explica que, antes de declarar la nulidad, el tribunal buscará interpretar el contrato a partir del lenguaje utilizado: \textit{the language of the document} —es decir, la redacción, la forma en que está escrito. El docente lo ejemplifica: si un contrato no lleva un nombre expreso, pero designa a las partes como distribuido y distribuidor, puede entenderse igualmente que se trata de un contrato de distribución, atendiendo a la letra del documento.

\begin{formula}{The four corners of the contract}
\textit{If ambiguity remains after looking within the four corners of the contract, such considerations may include the past dealings between the parties [y] relevant industry standards} —si la ambigüedad persiste luego de analizar el texto del contrato (\textit{the four corners}: literalmente, ``las cuatro esquinas'', es decir, todo el documento), el tribunal puede recurrir a las tratativas previas entre las partes (intercambio de correos, idas y vueltas de negociación) y a los usos y costumbres habituales del rubro de que se trate (\textit{industry standards}) —por ejemplo, el manual del franquiciante en un contrato de franquicia.
\end{formula}

\subsection{Validez versus exigibilidad (\textit{enforceability})}

El docente distingue dos planos que suelen confundirse: la validez del contrato y su exigibilidad. Un contrato puede ser válido —contar con todos sus elementos y estar bien celebrado— y, sin embargo, resultar \textit{unenforceable}: no exigible o no ejecutable ante un tribunal.

\begin{example}{Compraventa inmobiliaria no inscripta}
Se menciona el caso de una compraventa de un inmueble que no fue debidamente inscripta en el Registro de la Propiedad: el contrato es válido entre las partes, pero carece de las condiciones formales de oponibilidad frente a terceros. La cuestión de la exigibilidad se plantea recién cuando alguna de las partes lo lleva ante un tribunal.
\end{example}

El docente adelanta que el texto continúa con el desarrollo de la \textit{enforceability of contracts}: \textit{not all contracts, even if properly formed, will necessarily be enforced by the court} —no todos los contratos, aun teniendo todos sus elementos y siendo válidos, resultan finalmente exigibles ante un tribunal. Ese desarrollo —los motivos por los cuales un contrato válido puede no ser exigible— queda anunciado para la clase siguiente.

% ----------------------------------------------------------------------
\section*{Síntesis del desarrollo teórico}
\addcontentsline{toc}{section}{Síntesis del desarrollo teórico}

A lo largo de la clase, el docente recorre los cuatro elementos que el texto identifica para la formación de un contrato en el \textit{common law}: \textit{intention} (con sus presunciones diferenciadas según se trate de vínculos familiares o comerciales), \textit{offer and acceptance} (que juntos conforman el \textit{meeting of the minds}), \textit{consideration} (el intercambio de prestaciones de valor) y \textit{certainty of terms} (la exigencia de claridad y completitud de las cláusulas). En cada tramo remarca las diferencias con la estructura de la teoría general del contrato del derecho argentino —que no divide estos elementos del mismo modo— y anticipa que la clase siguiente abordará la \textit{enforceability of contracts}: los motivos por los cuales un contrato válidamente formado puede, aun así, no ser exigible ante un tribunal.

% ----------------------------------------------------------------------
\section*{Cuadro Cornell: repaso general de la unidad}
\addcontentsline{toc}{section}{Cuadro Cornell: repaso general de la unidad}

Formato \textit{Cornell Notes}: la columna izquierda plantea la pregunta o el concepto clave de cada tramo de la clase; la columna derecha desarrolla la respuesta o explicación correspondiente.

\begin{longtable}{
    >{\raggedright\arraybackslash\bfseries}p{0.32\textwidth}
    >{\raggedright\arraybackslash}p{0.58\textwidth}
}
\toprule
\textbf{Preguntas / palabras clave} & \textbf{Notas / respuestas} \\
\midrule
\endfirsthead

\toprule
\textbf{Preguntas / palabras clave} & \textbf{Notas / respuestas} \\
\midrule
\endhead

\midrule
\multicolumn{2}{r}{\textit{Continúa en la página siguiente}} \\
\endfoot

\bottomrule
\endlastfoot

¿Qué es la \textit{intention} en la formación del contrato? &
La intención real de las partes de celebrar un contrato vinculante (\textit{binding}). Una mera propuesta o idea vaga no alcanza; ante la duda, el tribunal evalúa pruebas (\textit{evidence}) y la naturaleza del vínculo entre las partes. \\

¿Cómo se presume la intención en vínculos familiares o afectivos? &
Como regla general, se presume que NO hay intención de contratar entre cónyuges, amigos cercanos o miembros de una familia. La presunción puede refutarse con pruebas (acuerdo escrito, intervención de abogados, actos que demuestren confianza en el acuerdo). \\

¿Qué resolvió el caso \textit{Merritt v. Merritt} (1970)? &
Pese al vínculo conyugal, la señora Merritt probó que existía un acuerdo escrito por el cual, al pagar íntegramente la hipoteca de la vivienda, el marido debía transferirle la propiedad. El tribunal reconoció la existencia de un contrato válido y ordenó la transferencia de dominio. \\

¿Cómo se presume la intención en vínculos comerciales? &
De manera opuesta a los vínculos familiares: en un entorno comercial (\textit{business setting}) se presume que SÍ existe intención de contratar. Las negociaciones previas (\textit{bargains}) —correos, presupuestos, señas— se toman como indicio de esa intención. \\

¿Qué es el \textit{meeting of the minds}? &
El entendimiento mutuo entre las partes sobre los términos y condiciones (\textit{terms and conditions}, es decir, las cláusulas) del contrato. Se logra mediante la concurrencia de una oferta (\textit{offer}) y una aceptación (\textit{acceptance}). \\

¿Qué requisitos debe reunir una oferta válida (\textit{offer})? &
Debe ser \textit{purposeful} (con intencionalidad de contratar) y expresar con claridad todos los términos del contrato (monto, plazo, objeto). Una propuesta imprecisa no alcanza a constituir una oferta en sentido estricto. \\

¿Qué significan los sufijos \textit{-or} y \textit{-ee} en los términos jurídicos ingleses? &
\textit{-or} identifica al sujeto activo que origina el negocio jurídico (\textit{offeror, employer, franchisor, grantor}); \textit{-ee} identifica al sujeto que recibe o es beneficiario (\textit{offeree, employee, franchisee}). \\

¿Cuándo caduca una oferta y qué es una \textit{counter offer}? &
Las ofertas pueden caducar tras un período de tiempo razonable, aunque no se haya fijado plazo expreso (ejemplo del collar reclamado 50 años después). Si la aceptación modifica los términos originales, no hay aceptación sino una contraoferta (\textit{counter offer}), que invierte los roles de \textit{offeror} y \textit{offeree}. \\

¿Cómo debe comunicarse la aceptación? &
De manera expresa (\textit{communicated expressly}), sea en forma oral o escrita, o de manera implícita mediante conductas o signos inequívocos; puede realizarla también un representante que actúe en nombre de la parte (\textit{on that person's behalf}). \\

¿Qué es \textit{consideration} y por qué no debe traducirse como ``consideración''? &
Es la contraprestación: el intercambio de promesas o prestaciones recíprocas entre las partes (\textit{each party promises something to the other}). ``Consideración'' es un falso amigo del inglés \textit{consideration}. \\

¿Qué significa que la \textit{consideration} deba ser ``\textit{something of value}''? &
Debe tratarse de algo susceptible de apreciación pecuniaria —un beneficio para quien lo recibe y, correlativamente, un \textit{detriment} (pérdida o desventaja en sentido amplio) para quien lo entrega—; un valor simbólico mínimo (como un caramelo) no alcanza. \\

¿La contraprestación siempre consiste en dinero? &
No. Puede consistir en un intercambio de favores, de bienes u objetos (\textit{items}), o en la abstención de realizar una determinada conducta (\textit{to refrain from doing something}), como en el caso del compromiso de no ocupar una parcela de un cementerio familiar. \\

¿Qué es la \textit{certainty of terms} y por qué es necesaria? &
La exigencia de que los términos del contrato sean claros y completos, de modo que ambas partes comprendan exactamente qué están contratando; la ambigüedad o vaguedad (\textit{ambiguity/vagueness}) en información clave puede volver el contrato incompleto o inválido. \\

¿Qué diferencia hay entre nulidad total y nulidad parcial (\textit{the contract will fail in whole or in part})? &
Si el vicio afecta solo una cláusula o estipulación, esa parte queda sin efecto (\textit{unenforceable}) y el resto del contrato subsiste; si el vicio es sustancial, el contrato entero puede resultar nulo (\textit{invalid}). \\

¿Con qué elementos cuenta el tribunal para interpretar un contrato ambiguo? &
Primero, con el lenguaje del propio documento (\textit{the four corners of the contract}); si la ambigüedad persiste, puede recurrir a las tratativas previas entre las partes (\textit{past dealings}) y a los usos y costumbres del rubro (\textit{relevant industry standards}). \\

¿Puede un contrato ser válido pero no exigible (\textit{unenforceable})? &
Sí. Un contrato puede reunir todos sus elementos y estar bien celebrado (válido) pero carecer de condiciones formales que lo hagan oponible o ejecutable ante un tribunal —por ejemplo, una compraventa inmobiliaria no inscripta registralmente frente a terceros. \\

\midrule
\multicolumn{2}{p{0.90\textwidth}}{
\textbf{Resumen:} La formación del contrato en el \textit{common law} se estructura sobre cuatro elementos que operan como eslabones de una misma cadena: la intención de contratar, sujeta a presunciones opuestas según el vínculo sea familiar/afectivo o comercial y siempre refutable mediante prueba; la oferta y la aceptación, que deben coincidir exactamente en sus términos para producir el \textit{meeting of the minds}, bajo pena de configurar una contraoferta; la \textit{consideration} o contraprestación, entendida como el intercambio de algo de valor —dinero, bienes, servicios o incluso una abstención— entre las partes; y la certeza de los términos, que exige claridad y completitud en las cláusulas, ya que el tribunal no redactará ni completará lo que las partes omitieron, sino que declarará la nulidad total o parcial del contrato. A lo largo de todo el desarrollo, el docente contrasta esta estructura con la teoría general del contrato del derecho argentino, que organiza estos mismos problemas de manera distinta (consentimiento, capacidad, objeto y causa), y remarca la importancia de identificar los falsos amigos del inglés jurídico —como \textit{consideration} o \textit{evidence}— para una lectocomprensión jurídica precisa. Queda planteado para la clase siguiente el desarrollo de la \textit{enforceability of contracts}: por qué un contrato válidamente formado puede, aun así, no ser exigible ante un tribunal.
} \\

\end{longtable}

\end{document}
